\section{Conclusion}
\label{sec:conclusion}

This dissertation evaluated whether a Docker-based NDT can enforce and measure controlled network conditions while remaining reproducible and extensible. For RQ1, repeated bandwidth, delay, and loss experiments produced distinct measured regimes with retained per-run evidence and uncertainty estimates; the results also showed that configuration values must not be conflated with end-to-end TCP, RTT, or loss observations. For RQ2, the same pipeline extended from direct and dual-router topologies to Germany50-derived paths and a complete 50-node/88-link instantiation. Routing claims remain bounded: 4,224 entries were dry-run validated, while real full-topology traffic covered only three representative paths. For RQ3, validated Qwen-compatible generation and 80 real-Docker control episodes demonstrated shared automation hooks, while the official OpenAI call failed with HTTP 429 and the threshold heuristic materially outperformed Q-learning.

The main conclusion is therefore conditional rather than universal. Within the evaluated configurations, the proposed system provides a reproducible workflow for controlled Docker network experimentation, validated routing, structured evidence, and guarded AI/control extensions. Its limitations---shared-host fidelity, small cohorts, representative Germany50 coverage, provider availability, missing Dashboard validation logs, and Q-learning underperformance---define the boundary of the contribution and the most credible directions for future work.
